\scanpage{48}\mspage{9}
construire les $\vec{\unsure{Z|\mathcal Z|\mathfrak Z}}_{s,\deleted{\unsure{\Delta}}}$ \inserted{$(s\in S)$}, et l'application
surjective
\[\vec I_{\deleted{\unsure{s}}}\setminus\vec I_s\longrightarrow\vec Z_{\unsure{s,|s;|s}}\tag*{(31)}\]
de l'ens des \unsure{pseudo-dr.|ps. dr.} orientées ne passant pas
par $s$, dans $\vec Z_s$, application compatible avec
l'action du groupe des signes $\pm1$.

Soit maintenant $\vec a$ un arc de
$K$, joignant $s$ à $s'$, on veut définir
\[\psi_{\vec a}:\vec Z_s\longrightarrow\vec Z_{s'}\]
\deleted{C'est ici qu'on a besoin de} \inserted{Supposons d'abord} l'existence
d'au moins une ps-droite qui ne passe
pas par $s$ ni $s'$. Considérons
\diagram{\begin{tikzcd}[ampersand replacement=\&,column sep=.9em,row sep=1.7em]
\& \& {\vec J=\vec I\setminus(\vec I_s\cup\vec I_{s'})}\arrow[dl]\arrow[dr] \& \\
{(32)} \& \vec Z_s\arrow[rr,dashed,"\psi_{\vec a}"] \& \& \vec Z_{s'}
\end{tikzcd}}
on obtient deux quotients de $\vec J$ — on
devra faire condition que c'est le $\overline{\textrm m}$, i.e.

(A4) si $\vec\Delta,\vec{\Delta'}$ sont deux ps. droites orientées
ne passant pas par les extrémités $s,s'$ de
l'arête $a$, qui induisent la $\overline{\textrm m}$ orientation en
$s$, elles induisent aussi la $\overline{\textrm m}$ orientation en $s'$.
% archive sheet number: 47
